\documentclass[10pt,a4paper]{amsart}
\usepackage{amsmath,amssymb,amsthm,graphicx}
% Retained minimal preamble for the Stacks Project chapter more-algebra.tex (More on Algebra).
% The chapter file pinned-original-source.tex carries only \input{preamble} on its first line, so no
% preamble text lives inside the chapter itself.  The Stacks Project uses one single shared preamble for
% every chapter; the snapshot stages it identically as sources/stacks-chow/preamble.tex and
% sources/stacks-etale-cohomology/preamble.tex.  The block below is copied mechanically, line for line,
% from sources/stacks-chow/preamble.tex lines 14-18 (the Xy-pic setup) and 154-217 (the native theorem
% declarations together with the mathematical macros).  Nothing was added, removed or redefined; the
% wrapper class and the amsmath/amssymb/amsthm/graphicx base packages are supplied by the assembler.
% Chapter: more-algebra.tex -- More on Algebra.
% For commutative diagrams we use Xy-pic
\usepackage[all]{xy}

% We use 2cell for 2-commutative diagrams.
\xyoption{2cell}

% Theorem environments.
%
\theoremstyle{plain}
\newtheorem{theorem}[subsection]{Theorem}
\newtheorem{proposition}[subsection]{Proposition}
\newtheorem{lemma}[subsection]{Lemma}

\theoremstyle{definition}
\newtheorem{definition}[subsection]{Definition}
\newtheorem{example}[subsection]{Example}
\newtheorem{exercise}[subsection]{Exercise}
\newtheorem{situation}[subsection]{Situation}

\theoremstyle{remark}
\newtheorem{remark}[subsection]{Remark}
\newtheorem{remarks}[subsection]{Remarks}

\numberwithin{equation}{subsection}

% Macros
%
\def\lim{\mathop{\mathrm{lim}}\nolimits}
\def\colim{\mathop{\mathrm{colim}}\nolimits}
\def\Spec{\mathop{\mathrm{Spec}}}
\def\Hom{\mathop{\mathrm{Hom}}\nolimits}
\def\Ext{\mathop{\mathrm{Ext}}\nolimits}
\def\SheafHom{\mathop{\mathcal{H}\!\mathit{om}}\nolimits}
\def\SheafExt{\mathop{\mathcal{E}\!\mathit{xt}}\nolimits}
\def\Sch{\mathit{Sch}}
\def\Mor{\mathop{\mathrm{Mor}}\nolimits}
\def\Ob{\mathop{\mathrm{Ob}}\nolimits}
\def\Sh{\mathop{\mathit{Sh}}\nolimits}
\def\NL{\mathop{N\!L}\nolimits}
\def\CH{\mathop{\mathrm{CH}}\nolimits}
\def\proetale{{pro\text{-}\acute{e}tale}}
\def\etale{{\acute{e}tale}}
\def\QCoh{\mathit{QCoh}}
\def\Ker{\mathop{\mathrm{Ker}}}
\def\Im{\mathop{\mathrm{Im}}}
\def\Coker{\mathop{\mathrm{Coker}}}
\def\Coim{\mathop{\mathrm{Coim}}}

% Boxtimes
%
\DeclareMathSymbol{\boxtimes}{\mathbin}{AMSa}{"02}

%
% Macros for moduli stacks/spaces
%
\def\QCohstack{\mathcal{QC}\!\mathit{oh}}
\def\Cohstack{\mathcal{C}\!\mathit{oh}}
\def\Spacesstack{\mathcal{S}\!\mathit{paces}}
\def\Quotfunctor{\mathrm{Quot}}
\def\Hilbfunctor{\mathrm{Hilb}}
\def\Curvesstack{\mathcal{C}\!\mathit{urves}}
\def\Polarizedstack{\mathcal{P}\!\mathit{olarized}}
\def\Complexesstack{\mathcal{C}\!\mathit{omplexes}}
% \Pic is the operator that assigns to X its picard group, usage \Pic(X)
% \Picardstack_{X/B} denotes the Picard stack of X over B
% \Picardfunctor_{X/B} denotes the Picard functor of X over B
\def\Pic{\mathop{\mathrm{Pic}}\nolimits}
\def\Picardstack{\mathcal{P}\!\mathit{ic}}
\def\Picardfunctor{\mathrm{Pic}}
\def\Deformationcategory{\mathcal{D}\!\mathit{ef}}
\usepackage{cleveref}
\crefname{originallemma}{original lemma}{original lemmas}
\Crefname{originallemma}{Original lemma}{Original lemmas}
\makeatletter
\newlabel{algebra-lemma-characterize-flat}{{00M5}{source}{}{sourcepointer-1}{}}
\newlabel{algebra-lemma-characterize-flat@cref}{{[originallemma][0][]00M5}{[1][0][]source}{}{}{}}
\makeatother
\crefname{originallemma}{original lemma}{original lemmas}
\Crefname{originallemma}{Original lemma}{Original lemmas}
\makeatletter
\newlabel{algebra-lemma-flat-base-change}{{00HI}{source}{}{sourcepointer-2}{}}
\newlabel{algebra-lemma-flat-base-change@cref}{{[originallemma][0][]00HI}{[1][0][]source}{}{}{}}
\makeatother
\crefname{originallemma}{original lemma}{original lemmas}
\Crefname{originallemma}{Original lemma}{Original lemmas}
\makeatletter
\newlabel{algebra-lemma-flat-tor-zero}{{00HL}{source}{}{sourcepointer-3}{}}
\newlabel{algebra-lemma-flat-tor-zero@cref}{{[originallemma][0][]00HL}{[1][0][]source}{}{}{}}
\makeatother
\crefname{originallemma}{original lemma}{original lemmas}
\Crefname{originallemma}{Original lemma}{Original lemmas}
\makeatletter
\newlabel{algebra-theorem-lazard}{{058G}{source}{}{sourcepointer-4}{}}
\newlabel{algebra-theorem-lazard@cref}{{[originallemma][0][]058G}{[1][0][]source}{}{}{}}
\makeatother
\crefname{originalremark}{original remark}{original remarks}
\Crefname{originalremark}{Original remark}{Original remarks}
\makeatletter
\newlabel{example-tor-change-rings}{{15192}{source}{}{sourcepointer-5}{}}
\newlabel{example-tor-change-rings@cref}{{[originalremark][0][]15192}{[1][0][]source}{}{}{}}
\makeatother
\crefname{originallemma}{original lemma}{original lemmas}
\Crefname{originallemma}{Original lemma}{Original lemmas}
\makeatletter
\newlabel{lemma-BL-faithful}{{25149}{source}{}{sourcepointer-6}{}}
\newlabel{lemma-BL-faithful@cref}{{[originallemma][0][]25149}{[1][0][]source}{}{}{}}
\makeatother
\crefname{originalremark}{original remark}{original remarks}
\Crefname{originalremark}{Original remark}{Original remarks}
\makeatletter
\newlabel{remark-what-you-get-for-general-modules}{{25676}{source}{}{sourcepointer-7}{}}
\newlabel{remark-what-you-get-for-general-modules@cref}{{[originalremark][0][]25676}{[1][0][]source}{}{}{}}
\makeatother
\title{Flatness of a module over a glueing pair is detected by its two localizations}
\author{The Stacks Project Authors}
\date{stacks-project:more-algebra.tex:local-snapshot-20261012; GNU FDL 1.2 or later}
\begin{document}
\maketitle
\setcounter{subsection}{1}
\noindent\textit{Editorial note.} All delivered mathematical text is the author's own, copied from the pinned Stacks Project chapter file more-algebra.tex; no statement or proof step was rewritten, reordered or removed.\par\medskip
\noindent\textit{Editorial note.} The chapter file pinned-original-source.tex carries only \textbackslash{}input\{preamble\} on its first line, so no preamble text lives inside the chapter. The retained file sources/stacks-more-algebra/retained-preamble.tex is a mechanical, line-for-line copy of the project's single shared preamble as staged in the snapshot; nothing was added, removed or redefined and no package was added.\par\medskip
\noindent\textit{Editorial note.} The source URL is the source-level Stacks Project URL, as required when more than one card is produced from one source; the per-claim tag is carried by claim\_key.\par\medskip
\noindent\textit{Editorial note.} The author's sentence 'the action on f on Tor\_i\textasciicircum{}R(...)' in the last line is reproduced unchanged. The seven references inside the counted proof span are recorded as non-load-bearing navigation pointers: algebra-lemma-flat-base-change, algebra-lemma-flat-tor-zero, algebra-lemma-characterize-flat, algebra-theorem-lazard and example-tor-change-rings resolve to the project's cross-chapter tag list, while remark-what-you-get-for-general-modules and lemma-BL-faithful are same-chapter labels resolved to their \textbackslash{}label lines in the pinned source (25676 and 25149).\par\medskip
\noindent\textit{Allowed standard prerequisites (editorial).}\begin{itemize}
\item Glueing pairs (R -> R', f), glueable modules and the glued module M-tilde. fixed earlier in the same section, used at proof lines 25709-25710 and 25720. Supplies the reduction from an arbitrary M to a glueable one, which is the first half of the converse direction.
\item The flatness criterion via Tor\_1\textasciicircum{}R(M, N) = 0, and Lazard's theorem that a flat module is a filtered colimit of finite free modules. used at proof lines 25729-25730 and 25751-25752. Converts the problem into a Tor-vanishing statement and lets the f-action be transferred to the base change.
\item The change-of-rings spectral sequence for Tor and the exact sequence 0 -> R -> R' oplus R\_f -> (R')\_f -> 0. used at proof lines 25733-25750. Computes Tor\_i\textasciicircum{}R(M tensor\_R R', N) from Tor\_i\textasciicircum{}R(R', N) and produces the f-action statement.
\end{itemize}\medskip
\ifdefined\hypertarget\hypertarget{sourcepointer-1}{}\fi
\noindent\textit{Source reference.} Original lemma 00M5: focused-seeds/AGP/sources/stacks-snapshot-20261011/tags.txt, source lines 811--811. Cross-chapter navigation pointer into the Stacks Project; the standard fact it supplies is applied inside the counted proof span and the argument resting on it stays there.\par\medskip
\ifdefined\hypertarget\hypertarget{sourcepointer-2}{}\fi
\noindent\textit{Source reference.} Original lemma 00HI: focused-seeds/AGP/sources/stacks-snapshot-20261011/tags.txt, source lines 644--644. Cross-chapter navigation pointer into the Stacks Project; the standard fact it supplies is applied inside the counted proof span and the argument resting on it stays there.\par\medskip
\ifdefined\hypertarget\hypertarget{sourcepointer-3}{}\fi
\noindent\textit{Source reference.} Original lemma 00HL: focused-seeds/AGP/sources/stacks-snapshot-20261011/tags.txt, source lines 647--647. Cross-chapter navigation pointer into the Stacks Project; the standard fact it supplies is applied inside the counted proof span and the argument resting on it stays there.\par\medskip
\ifdefined\hypertarget\hypertarget{sourcepointer-4}{}\fi
\noindent\textit{Source reference.} Original lemma 058G: focused-seeds/AGP/sources/stacks-snapshot-20261011/tags.txt, source lines 6779--6779. Cross-chapter navigation pointer into the Stacks Project; the standard fact it supplies is applied inside the counted proof span and the argument resting on it stays there.\par\medskip
\ifdefined\hypertarget\hypertarget{sourcepointer-5}{}\fi
\noindent\textit{Source reference.} Original remark 15192: sources/stacks-more-algebra/pinned-original-source.tex, source lines 15192--15192. Same-chapter navigation pointer; the statement it names is standard for the section and the argument that consumes it stays inside the counted proof span.\par\medskip
\ifdefined\hypertarget\hypertarget{sourcepointer-6}{}\fi
\noindent\textit{Source reference.} Original lemma 25149: sources/stacks-more-algebra/pinned-original-source.tex, source lines 25149--25149. Same-chapter navigation pointer; the statement it names is standard for the section and the argument that consumes it stays inside the counted proof span.\par\medskip
\ifdefined\hypertarget\hypertarget{sourcepointer-7}{}\fi
\noindent\textit{Source reference.} Original remark 25676: sources/stacks-more-algebra/pinned-original-source.tex, source lines 25676--25676. Same-chapter navigation pointer; the statement it names is standard for the section and the argument that consumes it stays inside the counted proof span.\par\medskip
\begin{lemma}
\label{lemma-BL-flat}
Let $(R \to R', f)$ be a glueing pair. Let $M$ be an $R$-module
which is not necessarily glueable for $(R \to R', f)$. Then $M$
is flat over $R$ if and only if $M \otimes_R R'$ is flat over $R'$
and $M_f$ is flat over $R_f$.
\end{lemma}

\begin{proof}
One direction of the lemma follows from
Algebra, Lemma \ref{algebra-lemma-flat-base-change}.
For the other direction, assume $M \otimes_R R'$
is flat over $R'$ and $M_f$ is flat over $R_f$.
Let $\tilde M$ be as in Remark \ref{remark-what-you-get-for-general-modules}.
If $\tilde M$ is flat over $R$, then applying
Algebra, Lemma \ref{algebra-lemma-flat-tor-zero}
to the short exact sequence
$0 \to \Ker(M \to \tilde M) \to M \to \tilde M \to 0$
we find that $\Ker(M \to \tilde M) \otimes_R (R' \oplus R_f)$ is zero.
Hence $M = \tilde M$ by Lemma \ref{lemma-BL-faithful} and we conclude.
In other words, we may replace $M$ by $\tilde M$ and assume
$M$ is glueable for $(R \to R', f)$. Let $N$ be a second $R$-module.
It suffices to prove that $\text{Tor}_1^R(M, N) = 0$, see
Algebra, Lemma \ref{algebra-lemma-characterize-flat}.

\medskip\noindent
The long exact sequence of Tors associated to the
short exact sequence $0 \to R \to R' \oplus R_f \to (R')_f \to 0$
and $N$ gives an exact sequence
$$
0 \to \text{Tor}_1^R(R', N) \to \text{Tor}_1^R((R')_f, N)
$$
and isomorphisms
$\text{Tor}_i^R(R', N) = \text{Tor}_i^R((R')_f, N)$
for $i \geq 2$. Since
$\text{Tor}_i^R((R')_f, N) = \text{Tor}_i^R(R', N)_f$
we conclude that $f$ is a nonzerodivisor on $\text{Tor}_1^R(R', N)$
and invertible on $\text{Tor}_i^R(R', N)$ for $i \geq 2$.
Since $M \otimes_R R'$ is flat over $R'$ we have
$$
\text{Tor}_i^R(M \otimes_R R', N) =
(M \otimes_R R') \otimes_{R'} \text{Tor}_i^R(R', N)
$$
by the spectral sequence of Example \ref{example-tor-change-rings}.
Writing $M \otimes_R R'$ as a filtered colimit of
finite free $R'$-modules (Algebra, Theorem \ref{algebra-theorem-lazard})
we conclude that $f$ is a nonzerodivisor
on $\text{Tor}_1^R(M \otimes_R R', N)$ and invertible on
$\text{Tor}_i^R(M \otimes_R R', N)$. Next, we consider
the exact sequence
$0 \to M \to M \otimes_R R' \oplus M_f \to M \otimes_R (R')_f \to 0$
coming from the fact that $M$ is glueable and the associated long exact
sequence of $\text{Tor}$. The relevant part is
$$
\xymatrix{
\text{Tor}_1^R(M, N) \ar[r] &
\text{Tor}_1^R(M \otimes_R R', N) \ar[r] &
\text{Tor}_1^R(M \otimes_R (R')_f, N) \\
& \text{Tor}_2^R(M \otimes_R R', N) \ar[r] &
\text{Tor}_2^R(M \otimes_R (R')_f, N) \ar[llu]
}
$$
We conclude that $\text{Tor}_1^R(M, N) = 0$ by our remarks above on the
action on $f$ on $\text{Tor}_i^R(M \otimes_R R', N)$.
\end{proof}

\end{document}
